\documentclass[11pt]{article}
\usepackage{preamble}
\usepackage{multicol}
\setlength{\parskip}{4pt}

\begin{document}

\daybanner{AP Statistics \textperiodcentered\ Unit 1 \textperiodcentered\ Topics 1.1--1.7 (CED 1.1--1.7) \textperiodcentered\ remediation sheet \textperiodcentered\ TEACHER KEY}{Screen Time, Two Deletions}

\sectionbanner{CED TETHER}
\begin{multicols}{2}
{\footnotesize\setlength{\parskip}{0pt}
\begin{itemize}[leftmargin=*,itemsep=0pt,topsep=0pt]
\item VAR-1.A / UNC-1.A: identify the question to be answered and the individuals and variables in a context.
\item VAR-1.B: classify variables as categorical or quantitative; a label or identifier is not a measurement.
\item UNC-1.H / 1.I: compare groups of unequal size with relative frequencies, not raw counts.
\item UNC-1.J: describe a distribution with shape, center, variability, and unusual features.
\item UNC-1.K: the median is resistant to unusual values; the mean is not, so report the center that fits the data.
\end{itemize}}
\end{multicols}
\clearpage
\small
\textbf{Essential question:} What does a number mean, and what changes when you delete the values you do not like?\par
\textbf{DOK-3 skill:} 4.B \textperiodcentered\ \textbf{FRQ pattern:} {\small\texttt{judge-\allowbreak{}a-\allowbreak{}trimmed-\allowbreak{}report-\allowbreak{}using-\allowbreak{}context-\allowbreak{}resistance-\allowbreak{}and-\allowbreak{}relative-\allowbreak{}frequency}}\par
\textbf{DOK rationale:} Part (d) requires a defended decision: whether a report that deleted two values is responsible, which statistic to report, and what the changed percent means for the principal's actual question. It coordinates context, individuals-versus-variables, relative frequency, and resistance rather than recalling any one of them.\par

\frameworkphaseheader{First take $\rightarrow$ rules + (a)--(d) $\rightarrow$ turn in}{1 $\rightarrow$ 3}{32}{%
\begin{itemize}[leftmargin=1.5em,itemsep=2pt,topsep=2pt]
  \item Hand out the sheet; read the scenario aloud once. Start the first take at the bell (ungraded).
  \item Point at the printed rules box and say only that every rule will be needed. Do not interpret the rules.
  \item Allow 25 minutes for (a)--(d). Circulate; answer questions with questions. Collect one sheet; score (d) only, E/P/I by a human.
\end{itemize}
}{%
\begin{itemize}[leftmargin=1.5em,itemsep=2pt,topsep=2pt]
  \item Commit one sentence before any discussion.
  \item Work (a) to (d) in order; each part supplies a piece the next part needs.
  \item Use the checklists and frames on the page; every blank is theirs to fill.
\end{itemize}
}{%
\begin{itemize}[leftmargin=1.5em,itemsep=2pt,topsep=2pt]
  \item Who exactly was measured, and what is one row of the data?
  \item Which of your variables could you add up, and which could you only count?
  \item If the two groups were the same size, would your comparison change?
  \item Which number moved a lot and which barely moved when the two values were deleted? Why?
\end{itemize}
}{Circulate and ask for evidence linked to the printed rules. This is the human-graded DOK-3 channel; do not send it to an AI grader or auto-score it.}

\begin{callout}{\IconWarn\ LOOK FOR}{calloutred}
\begin{itemize}[leftmargin=1.5em,itemsep=2pt,topsep=2pt]
  \item Writing a number with no who/what/units (the 30\% left floating).
  \item Calling the student ID or the device name a quantitative variable.
  \item Comparing 9 freshmen to 3 seniors by counts and declaring freshmen three times worse.
  \item Saying the mean is resistant, or that deleting the two values was fine because the median barely moved.
  \item Describing the distribution without mentioning the two unusual values at all.
\end{itemize}
\end{callout}

\sectionbanner{SCORING PART (d) --- E / P / I}

\textbf{Expected elements}
\begin{itemize}[leftmargin=1.5em,itemsep=2pt,topsep=2pt]
  \item \textbf{e1} (required) --- Gives the 26\% its context: who (students in the sample), what (daily screen time over 240 minutes), and the question it answers.
  \item \textbf{e2} (required) --- Chooses a center and justifies it with resistance, consistent with part (c).
  \item \textbf{e3} (required) --- Interprets the big mean change versus the small median change as the effect of the two unusual values.
  \item \textbf{e4} (required) --- States that the deletion must be disclosed and justified (logging errors) and that unusual values still belong in the description.
\end{itemize}

{\small\renewcommand{\arraystretch}{1.25}
\noindent\begin{tabular}{|p{0.5in}|p{5.9in}|}
\hline
\textbf{E} & All four elements, in the context\textquotesingle s words, with a clear decision. \\ \hline
\textbf{P} & A decision with two or three elements, or all four with the context or the disclosure vague. \\ \hline
\textbf{I} & A bare number or a decision without resistance, context, or the disclosure. \\ \hline
\end{tabular}}

\clearpage
\normalsize
\focusbanner{TODAY'S PROBLEM --- annotated}

A school nurse recorded, for a random sample of 40 students, three things: grade level (Freshman or Senior), the device they use most (Phone, Laptop, or Console), and daily screen time in minutes from the phone\textquotesingle s own report. The principal asked one question: \emph{Is too much screen time a concern for our students?} The nurse\textquotesingle s first report said only: \emph{12 of the 40 are over 240, that is 30\%, so yes.}\par\smallskip Group sizes: 25 freshmen, 15 seniors. Over 240 minutes: 9 freshmen, 3 seniors.\par\smallskip The two largest values in the histogram (a 900-minute and a 720-minute entry) came from two students who logged a whole weekend by mistake. The nurse deleted both and re-computed: the mean fell from 205 to 178 minutes; the median moved from 176 to 174 minutes; the headline became \emph{10 of 38 are over 240, that is 26\%.}\par
\begin{center}
\begin{tikzpicture}
\begin{axis}[width=0.96\linewidth,height=1.45in,ybar interval,x tick label as interval=false,bar shift=0pt,xmin=0,xmax=960,
 ymin=0,ymax=16,xtick={0,120,240,360,480,600,720,840,960},ytick={0,4,8,12,16},xlabel={Daily screen time (minutes)},ylabel={Students},
 tick label style={font=\scriptsize},label style={font=\scriptsize},title={All 40 students, before any deletion},title style={font=\small}]
\addplot[fill=calloutblue,draw=dayblue] coordinates {(0,3) (120,14) (240,12) (360,6) (480,3) (600,0) (720,1) (840,1) (960,0)};
\end{axis}\end{tikzpicture}
\end{center}
\begin{firsttakebox}{FIRST TAKE (5 min)}
Before any calculation: does the nurse\textquotesingle s first sentence answer the principal\textquotesingle s question? Name one thing the sentence leaves out.\par
\answer{Ungraded commitment; accept any position that names something the first sentence leaves out.}
\end{firsttakebox}

\textit{Rules callout ``RULES YOU MUST USE (NO ANSWERS HERE)'' is printed on the student sheet.}\par\medskip

\dokbadge{1}~\textbf{(a)} Fill the table: who or what are the individuals, and for each recorded variable, what ONE value looks like and whether it is categorical or quantitative. Checklist: an individual is one thing measured; a name, label, or ID is not a measurement; quantitative values could sensibly be averaged.\par\smallskip \begin{tabular}{|p{1.3in}|p{2.2in}|p{1.6in}|}\hline \textbf{Individuals:} & \multicolumn{2}{l|}{\hrulefill} \\ \hline \textbf{Variable} & \textbf{One value looks like} & \textbf{Categorical / Quantitative} \\ \hline Grade level & & \\[10pt] \hline Device used most & & \\[10pt] \hline Daily screen time & & \\[10pt] \hline \end{tabular}\par
\answer{Individuals: the 40 students in the sample. Grade level: one value is Freshman or Senior, categorical. Device used most: one value is Phone, Laptop, or Console, categorical. Daily screen time: one value is a number of minutes such as 176, quantitative (it can be averaged).}

\dokbadge{2}~\textbf{(b)} A senior says: \emph{Nine freshmen are over 240 and only three seniors are, so freshmen are three times as bad.} Fix the comparison. Frame: \emph{Of the \underline{\hspace{0.4in}} freshmen, \underline{\hspace{0.4in}} are over 240, which is \underline{\hspace{0.4in}}\%. Of the \underline{\hspace{0.4in}} seniors, \underline{\hspace{0.4in}} are over 240, which is \underline{\hspace{0.4in}}\%. So the fair comparison is \underline{\hspace{0.9in}} because the groups \underline{\hspace{0.9in}}.} Then explain in one sentence why the count 9 versus 3 can mislead here.\par
\answer{Of the 25 freshmen, 9 are over 240, which is 36\%. Of the 15 seniors, 3 are over 240, which is 20\%. The fair comparison is 36\% versus 20\% because the groups have different sizes; 9 versus 3 exaggerates the gap (three times as many students, not three times as bad), and freshmen are still higher but by less than the counts suggest.}

\dokbadge{2}~\textbf{(c)} Deleting the two values moved the mean by 27 minutes and the median by 2 minutes. Which of these two statistics is resistant to unusual values? Explain what in the histogram caused the big move, using the words \emph{resistant} and \emph{unusual values}. Do not say which one the nurse should report yet.\par
\answer{The median is resistant; the mean is not. The two unusual values at 720 and 900 minutes sit far to the right of the rest of the distribution, so they pull the mean upward. Removing them drops the mean by 27 minutes, while the median, which depends only on the middle of the ordered data, moves by 2.}

\dokbadge{3}~\textbf{(d)} $\star$ Judge the nurse\textquotesingle s trimmed report in three to five sentences. Before you write, check: (1) WHO, WHAT, and the UNITS behind the 26\%; (2) what question 26\% answers, and whether that is the principal\textquotesingle s question; (3) which center to report, tied to your answer in (c); (4) whether deleting the two mistaken values was responsible and what the report must SAY about them. Use the frame below if it helps.\par
\answer{Report the median (about 174--176 minutes) because it is resistant and describes the typical student either way. Removing the two values changed the mean a lot but barely changed the median, which shows the mean was being pulled by two unusual values, not by typical students. The 26\% means that 10 of the 38 remaining students in the sample reported more than 240 minutes per day. Deleting the two entries is defensible because they were logging errors, not real days, but the report must say that two values were removed and why, and it should still note that unusual values existed; only then does the number connect to the principal\textquotesingle s question about whether screen time is a concern.}

\begin{summaryexitbox}{EXIT}
One sentence I will add to my next written answer so a number is never left without context: \hrulefill
\end{summaryexitbox}

\end{document}
